\documentclass[11pt]{article}

\usepackage[T1]{fontenc}
\usepackage[utf8]{inputenc}
\usepackage{lmodern}
\usepackage[a4paper,margin=28mm]{geometry}
\usepackage{microtype}
\usepackage{amsmath,amssymb,amsthm,mathtools}
\usepackage{booktabs}
\usepackage{array}
\usepackage{enumitem}
\usepackage{xcolor}
\usepackage{hyperref}
\usepackage{fancyhdr}

\definecolor{linkblue}{RGB}{26,78,130}
\definecolor{statusred}{RGB}{135,25,35}
\hypersetup{
  colorlinks=true,
  linkcolor=linkblue,
  citecolor=linkblue,
  urlcolor=linkblue,
  pdftitle={Finite Spectral-Gap Certificates in a Riemann Prime-Resolvent Programme},
  pdfauthor={Lluis Eriksson}
}

\pagestyle{fancy}
\fancyhf{}
\fancyhead[L]{Finite spectral-gap certificates}
\fancyhead[R]{L. Eriksson}
\fancyfoot[C]{\thepage}
\fancypagestyle{plain}{
  \fancyhf{}
  \fancyhead[L]{Finite spectral-gap certificates}
  \fancyhead[R]{L. Eriksson}
  \fancyfoot[C]{\thepage}
}
\setlength{\headheight}{14pt}
\setlength{\emergencystretch}{3em}

\newtheorem{theorem}{Theorem}[section]
\newtheorem{proposition}[theorem]{Proposition}
\newtheorem{lemma}[theorem]{Lemma}
\newtheorem{corollary}[theorem]{Corollary}
\theoremstyle{definition}
\newtheorem{definition}[theorem]{Definition}
\newtheorem{remark}[theorem]{Remark}

\newcommand{\RR}{\mathbb{R}}
\newcommand{\CC}{\mathbb{C}}
\newcommand{\XiR}{\mathcal{S}_{\Xi}}
\newcommand{\ip}[2]{\left\langle #1,#2\right\rangle}
\newcommand{\norm}[1]{\left\lVert #1\right\rVert}
\newcommand{\defect}{\mathfrak d}

\title{\textbf{Finite Spectral-Gap Certificates in a\\
Riemann Prime-Resolvent Programme}\\[4pt]
\large Lean verification, a subsequential criterion, and a falsification campaign}
\author{Llu\'is Eriksson\\
\small Independent researcher\\
\small \url{https://github.com/lluiseriksson/riemann-prime-resolvent}}
\date{Version 1.0 -- 6 August 2026}

\begin{document}
\maketitle
\thispagestyle{empty}

\begin{center}
\fcolorbox{statusred}{white}{\parbox{0.91\textwidth}{
\centering\textbf{Scope statement.}
This article does not prove the Riemann hypothesis.  It proves finite
Hilbert-space and spectral-gap estimates, reports their Lean 4 formalization,
and records a conditional analytic route whose concrete asymptotic inputs
remain open.}}
\end{center}

\begin{abstract}
We isolate a finite-dimensional estimate used in a prime-resolvent research
programme related to the Riemann $\Xi$-function.  Let $T$ be a real symmetric
operator, let $g$ be a normalized ground eigenvector with eigenvalue $\lambda_0$,
and let $u$ be a normalized trial vector.  If the ground state is separated
from the rest of the spectrum by $\Delta>0$, then
\[
 \Delta\bigl(1-\ip{g}{u}^{2}\bigr)
 \leq \ip{Tu}{u}-\lambda_0,
 \qquad
 \norm{g-u}\leq
 \sqrt{\frac{2(\ip{Tu}{u}-\lambda_0)}{\Delta}}
\]
after sign alignment.  Returning from a normalized Galerkin projection to a
complete trial vector costs at most twice the projection tail.  We formalize
these statements in Lean 4, including a guard proving that a positive indexed
gap forces a one-dimensional ground eigenspace.  Since Mathlib orders the
finite eigenvalue list decreasingly, a single inequality between the final two
entries suffices to certify the universal gap.

We then state a repaired subsequential criterion: along a cofinal family of
finite approximants, a defect
\[
 \defect_j=\norm{k_j}_2\sqrt{2r_j/\Delta_j}+2\,\mathrm{tail}_j
\]
with suitable polynomial decay, together with finite parity, non-vanishing and
local analytic convergence hypotheses, would imply the Riemann hypothesis.
Those hypotheses are not established here.  Finally, we report a registered
finite falsification campaign.  Its global verdict was \textsc{inconclusive};
an independent polynomial-root witness diagnosed one missing scan root at
$(\lambda,N)=(3,40)$ as a base-resolution failure.  The numerical experiment
is not used in any proof.
\end{abstract}

\noindent\textbf{Keywords:} Riemann hypothesis; spectral gap; Rayleigh
quotient; Galerkin approximation; Stieltjes transform; Hausdorff moments;
formal verification; Lean 4.

\section{Introduction and claim boundary}

Write
\[
 \Xi(z)=\xi\!\left(\frac12+iz\right),
 \qquad
 \XiR(x)=\frac{1}{2\sqrt{x}}
 \frac{\xi'}{\xi}\!\left(\frac12+\sqrt{x}\right),
 \quad x>0,
\]
where $\xi$ is the completed Riemann zeta function.  A companion criterion
packages the real-zero question for $\Xi$ as a one-point Hausdorff moment
problem.  The construction side of the programme asks whether positive finite
spectral approximants can converge non-circularly to $\XiR$.

The purpose of this paper is narrower.  It closes a finite spectral estimate
that had previously appeared only as an informal obligation.  It also records
exactly what the estimate does not provide.  In particular, it does not
construct the required infinite operator family, prove the necessary
prime-archimedean cancellation, or establish convergence to $\Xi$.

The contributions are:
\begin{enumerate}[label=(\roman*)]
  \item an elementary finite spectral-gap inequality and its Galerkin-tail
  completion;
  \item a formal degeneracy guard: the positive indexed gap implies that the
  ground eigenspace has dimension one;
  \item a reduction of the universal gap certificate to the last adjacent gap
  in Mathlib's decreasingly ordered eigenvalue list;
  \item a Lean 4 implementation checked without project-specific axioms or
  placeholders;
  \item a precise statement of a conditional subsequential route and a
  separately labelled numerical falsification campaign.
\end{enumerate}

Classical background on $\zeta$ and $\xi$ may be found in
\cite{riemann1859,titchmarsh1986}.  The Hausdorff moment theorem is due to
Hausdorff \cite{hausdorff1921}; standard references include
\cite{shohat1943,akhiezer1965}.  The finite spectral estimates used below are
in the family of standard Rayleigh, Temple and Kato arguments
\cite{kato1995}.

\section{The one-point criterion and the open construction arrow}

Fix $x_0>0$ and define the normalized derivative jet
\[
 b_n(x_0)=x_0^n\frac{(-1)^n}{n!}\XiR^{(n)}(x_0).
\]
For a sequence $b$, define $D^0b_n=b_n$ and
\[
 D^{k+1}b_n=D^kb_n-D^kb_{n+1}.
\]
The documented analytic argument in the companion criterion layer has the
following form.

\begin{proposition}[Documented one-point criterion]\label{prop:hausdorff}
For a fixed $x_0>0$, subject to the stated branch and logarithmic-derivative
conventions, the following are equivalent:
\begin{enumerate}[label=(\alph*)]
  \item all nontrivial zeros of $\zeta$ lie on the critical line;
  \item $(b_n(x_0))_{n\geq0}$ is a Hausdorff moment sequence on $[0,1]$;
  \item $D^kb_n(x_0)\geq0$ for all $n,k\geq0$.
\end{enumerate}
\end{proposition}

Under the Riemann hypothesis the representing atoms are obtained from the
ordinates $\gamma_j$ by
\[
 v_j=\frac{x_0}{x_0+\gamma_j^2}\in(0,1),
 \qquad
 b_n(x_0)=\sum_j\frac{1}{x_0+\gamma_j^2}\,v_j^n.
\]
Conversely, a representing measure yields a Stieltjes transform holomorphic
off the negative real axis and agreeing locally with $\XiR$.

Proposition~\ref{prop:hausdorff} is a documented analytic theorem, not a fully
closed Lean theorem in the present repository.  The difficult construction
arrow is
\[
 \text{prime-built finite spectra}
 \longrightarrow
 \text{positive Stieltjes limit agreeing with }\XiR.
\]
The finite result below addresses one estimate inside that arrow.

\section{Finite spectral-gap control}

Let $E$ be a finite-dimensional real Hilbert space.  Let
$T:E\to E$ be symmetric and choose an orthonormal eigenbasis
$(e_i)_{i\in I}$ with eigenvalues $(\lambda_i)_{i\in I}$.  Fix an index
$i_0$ and write
\[
 g=e_{i_0},\qquad \lambda_0=\lambda_{i_0}.
\]
Let $u\in E$ satisfy $\norm{u}=1$ and put
\[
 c_i=\ip{e_i}{u},
 \qquad
 r=\ip{Tu}{u}-\lambda_0.
\]

\begin{theorem}[Finite Rayleigh-gap inequality]\label{thm:gap}
Suppose that for some $\Delta\in\RR$,
\[
 \lambda_0+\Delta\leq\lambda_i
 \quad\text{for every }i\neq i_0.
\]
Then
\[
 \Delta\bigl(1-\ip{g}{u}^2\bigr)\leq r.
\]
\end{theorem}

\begin{proof}
Parseval's identity and diagonalization give
\[
 \sum_i c_i^2=1,
 \qquad
 \ip{Tu}{u}=\sum_i\lambda_i c_i^2.
\]
Consequently
\begin{align*}
 r
 &=\sum_i(\lambda_i-\lambda_0)c_i^2\\
 &=\sum_{i\neq i_0}(\lambda_i-\lambda_0)c_i^2\\
 &\geq \Delta\sum_{i\neq i_0}c_i^2
 =\Delta(1-c_{i_0}^2).
\end{align*}
\end{proof}

\begin{corollary}[Distance bound]\label{cor:distance}
Assume the hypotheses of Theorem~\ref{thm:gap}, $\Delta>0$, and choose the
sign of $g$ so that $\ip{g}{u}\geq0$.  Then $r\geq0$ and
\[
 \norm{g-u}\leq\sqrt{\frac{2r}{\Delta}}.
\]
\end{corollary}

\begin{proof}
For $c=\ip{g}{u}\in[0,1]$,
\[
 \norm{g-u}^2=2(1-c)\leq2(1-c^2)\leq\frac{2r}{\Delta}.
\]
\end{proof}

\subsection{Why the gap hypothesis encodes simplicity}

The strict positivity of $\Delta$ is essential.  If the ground eigenvalue
were degenerate, a normalized $u$ could lie in the same eigenspace while being
orthogonal to a chosen $g$.  Then $r=0$ but $\norm{g-u}=\sqrt2$, contradicting
the claimed distance estimate.

The indexed gap rules out this counterexample because an eigenbasis enumerates
multiplicity.

\begin{corollary}[Degeneracy guard]\label{cor:simple}
If $\Delta>0$ and
$\lambda_0+\Delta\leq\lambda_i$ for every $i\neq i_0$, then $i_0$ is the
unique index with eigenvalue $\lambda_0$.  For a symmetric operator this
implies
\[
 \dim\ker(T-\lambda_0 I)=1.
\]
\end{corollary}

In the formal proof the final statement is not inferred from informal
language.  Mathlib identifies the cardinality of the filtered eigenvalue list
with the finite dimension of the eigenspace.  The positive gap makes that
filtered list the singleton $\{i_0\}$.

\subsection{One adjacent inequality is sufficient}

Mathlib's finite spectral theorem orders the eigenvalues of a symmetric
operator decreasingly.  Thus, for a list
\[
 \lambda_0^{\downarrow}\geq\lambda_1^{\downarrow}\geq\cdots
 \geq\lambda_{m-2}^{\downarrow}\geq\lambda_{m-1}^{\downarrow},
\]
the ground eigenvalue is the last entry.

\begin{theorem}[Adjacent-gap reduction]\label{thm:adjacent}
Let $m\geq2$ and suppose the finite eigenvalue list is antitone.  If
\[
 \lambda_{m-1}^{\downarrow}+\Delta
 \leq\lambda_{m-2}^{\downarrow},
\]
then for every $i\neq m-1$,
\[
 \lambda_{m-1}^{\downarrow}+\Delta
 \leq\lambda_i^{\downarrow}.
\]
If $\Delta>0$, the last eigenvalue is simple.
\end{theorem}

\begin{proof}
Every index other than the last is at most the penultimate index.  Antitonicity
therefore gives
$\lambda_{m-2}^{\downarrow}\leq\lambda_i^{\downarrow}$, and the result follows
by transitivity.  Simplicity follows from Corollary~\ref{cor:simple}.
\end{proof}

This reduction matters computationally: an interval certificate for the two
lowest eigenvalues can supply the entire indexed-gap premise.

\section{Galerkin projection and the exact defect}

Let $k$ be a complete trial vector and $P_Nk$ its finite Galerkin projection.
Assume $P_Nk\neq0$ and set
\[
 u=\frac{P_Nk}{\norm{P_Nk}},
 \qquad
 \mathrm{tail}=\norm{k-P_Nk}.
\]
Let $g_N$ be the aligned normalized ground eigenvector of the finite symmetric
operator, let $\Delta_N>0$ be its spectral gap, and let
\[
 r_N=\ip{T_Nu}{u}-\lambda_{0,N}.
\]

\begin{theorem}[Finite Galerkin defect]\label{thm:galerkin}
Under the preceding assumptions,
\[
 \norm{\norm{k}g_N-k}
 \leq
 \norm{k}\sqrt{\frac{2r_N}{\Delta_N}}
 +2\,\mathrm{tail}.
\]
The same positive adjacent-gap certificate proves simultaneously that the
finite ground eigenspace is one-dimensional.
\end{theorem}

\begin{proof}
Corollary~\ref{cor:distance} gives
\[
 \norm{g_N-u}\leq\sqrt{2r_N/\Delta_N}.
\]
Insert $\norm{k}u$, $P_Nk=\norm{P_Nk}u$, and $k$ into a triangle inequality:
\begin{align*}
 \norm{\norm{k}g_N-k}
 &\leq \norm{k}\norm{g_N-u}
 +\bigl|\norm{k}-\norm{P_Nk}\bigr|
 +\norm{P_Nk-k}\\
 &\leq \norm{k}\sqrt{\frac{2r_N}{\Delta_N}}
 +2\,\norm{k-P_Nk}.
\end{align*}
The second inequality uses the reverse triangle inequality.  No small-tail
assumption such as $\mathrm{tail}\leq\norm{k}/2$ is needed.
\end{proof}

We therefore define the finite defect
\begin{equation}\label{eq:defect}
 \defect_N=\norm{k}\sqrt{\frac{2r_N}{\Delta_N}}
 +2\,\mathrm{tail}.
\end{equation}
The factor $2$ is not decorative: one tail term changes the scale
$\norm{P_Nk}$ to $\norm{k}$ and the other returns from $P_Nk$ to $k$.

\section{A conditional subsequential route}

The finite inequality is useful only if it is coupled to a concrete family.
We record the current conditional route because it identifies the exact
asymptotic obligation rather than hiding it in a generic convergence premise.

Let $(\lambda_j,N_j)$ be a cofinal diagonal of finite approximants.  For the
$j$th projected normalized trial vector define $r_j$, $\Delta_j$ and
$\mathrm{tail}_j$ as above, and put
\[
 \defect_j=\norm{k_j}_2\sqrt{\frac{2r_j}{\Delta_j}}
 +2\,\mathrm{tail}_j.
\]

\begin{proposition}[Conditional $\mathrm{RATE}_{\Sigma}$ criterion]
Suppose there are $0<a_1<a_2<a_3<1/2$ such that:
\begin{enumerate}[label=(P\arabic*)]
  \item along the cofinal diagonal, the relevant finite ground states are
  simple, have the required parity, and satisfy the non-vanishing condition
  needed by the finite zero-identification theorem;
  \item the transforms $\widehat{k}_j$ converge uniformly to $\Xi$ on a compact
  two-dimensional neighborhood of the segment $i[a_1,a_2]$ contained in
  $|\operatorname{Im}z|\leq a_3$;
  \item the same approximants satisfy
  \[
   \lambda_j^{a_3}\sqrt{\log\lambda_j}\,\defect_j\longrightarrow0.
  \]
\end{enumerate}
Then the documented Stieltjes and identity-theorem argument implies the
Riemann hypothesis.
\end{proposition}

This proposition is not presently a kernel-checked end-to-end theorem.  The
finite estimate~\eqref{eq:defect} is formalized, but the parity and
non-vanishing family, the local convergence to $\Xi$, the infinite atomic
compactness instantiation, and the decay in (P3) remain open.  In particular,
the proposition must not be quoted as an unconditional theorem about a
concrete operator.

The subsequential form reduces the verification surface: it uses finite real
zeros directly and requires decay only on a cofinal diagonal.  It does not
remove the arithmetic content.  Any positive power decay in (P3) still
requires signed cancellation not supplied by positive majorants alone.

\section{Lean 4 formalization}

The results of Sections 3 and 4 are implemented in the repository's
\path{SpectralDefect.lean} module.  The implementation uses
Mathlib's finite spectral theorem, orthonormal eigenvector basis, Parseval
identity and eigenvalue multiplicity count.  Lean and Mathlib are described in
\cite{demoura2015,mathlib2020}.

The main declaration chain is:
\begin{center}
\begin{tabular}{>{\raggedright\arraybackslash}p{0.43\textwidth}
                  >{\raggedright\arraybackslash}p{0.49\textwidth}}
\toprule
Declaration & Mathematical content\\
\midrule
\path{spectralGap_mul_one_sub_inner_sq_le_rayleighExcess_of_eigenbasis}
& Theorem~\ref{thm:gap} in an arbitrary orthonormal eigenbasis.\\[3pt]
\path{spectralGap_mul_one_sub_inner_sq_le_rayleighExcess_of_isSymmetric}
& Instantiation with Mathlib's finite spectral theorem.\\[3pt]
\path{finrank_groundEigenspace_eq_one_of_positive_gap}
& The ground eigenspace has dimension one.\\[3pt]
\path{antitone_gap_from_penultimate}
& The adjacent-gap reduction of Theorem~\ref{thm:adjacent}.\\[3pt]
\path{norm_scaled_eigenvectorBasis_sub_complete_le_rayleighGapDefect}
& The finite Galerkin estimate~\eqref{eq:defect}.\\[3pt]
\path{simpleGround_and_norm_scaled_eigenvectorBasis_sub_complete_le_of_adjacentGap}
& Combined adjacent-gap, simplicity and Galerkin certificate.\\
\bottomrule
\end{tabular}
\end{center}

The project oracle prints the axioms of each public theorem.  For the
declarations above it reports only the standard logical dependencies
\path{propext}, \path{Classical.choice} and \path{Quot.sound}.  The checked
source contains no \path{sorry}, \path{admit} or project-specific axiom.

This verification claim concerns the finite theorems in this article.  It does
not certify Proposition~\ref{prop:hausdorff}, the conditional
$\mathrm{RATE}_{\Sigma}$ criterion as an end-to-end result, or the missing
concrete operator construction.

\section{Registered finite falsification campaign}

A separate campaign tested a concrete finite spectral arm.  Its role was
falsification: predetermined gates could kill that arm at the probed scales,
but no finite outcome could prove the Riemann hypothesis or the asymptotic
criterion.

The registered grid was
\[
 \lambda\in\{2,\sqrt8,3,\sqrt{13}\},
 \qquad
 N\in\{20,40,80,120\},
\]
with a base and a high-resolution scan for each of the sixteen grid cells.
The global output was
\begin{quote}\small\ttfamily
T3-VERDICT: INCONCLUSIVE conv=S-half;x0=1;\\
grid=L\{2,sqrt8,3,sqrt13\}xN\{20,40,80,120\}\\
detail=\{"invalid": ["3:40", "3:80", "3:120"], "missing": []\}
\end{quote}
The three invalid entries were base scans at $\lambda=3$ with a deficit of
one detected positive root.  High-resolution replicas were valid.  Across the
thirteen valid cells, the registered kill observables reported no low-energy
spurious modes and zero low-energy mass under the specified threshold.  This
is a non-kill result at finite scales, not evidence of the required limit.

\subsection{Polynomial completeness witness at $(3,40)$}

The A4 witness replaced sign-change scanning by the characteristic polynomial
of the finite problem.  For $(\lambda,N)=(3,40)$, the reported certificate had
degree $80$ and found $80$ real roots, split into $40$ positive and $40$
negative roots, with $P(0)\neq0$.  The root missed by the base scan was
\[
 z=82.9103880394\ldots,
\]
whose nearest neighbor was approximately $1.8251$ away.  The high-resolution
scan found it; the base scan did not.  The mechanical A4 classification was
\textsc{base-resolution}, not a de-duplication merge, window failure or even
multiplicity.

This diagnosis is certified for $(3,40)$.  The matching deficit pattern at
$(3,80)$ and $(3,120)$ is suggestive but is not promoted here to a certified
cause without separate completeness witnesses.  Moreover, A4 was explicitly
non-gating: it did not alter the registered \textsc{inconclusive} verdict.

\section{What remains open}

The finite theorem removes one informal step but leaves the decisive
mathematics intact.  A complete argument would still require:
\begin{enumerate}[label=(\arabic*)]
  \item a concrete Hilbert space and operator family with domains and
  self-adjointness proved non-circularly;
  \item parity, non-vanishing and adjacent positive-gap certificates along an
  infinite cofinal diagonal;
  \item certified control of $r_j$, $\Delta_j$ and $\mathrm{tail}_j$ strong
  enough for $\mathrm{RATE}_{\Sigma}$;
  \item convergence of the transforms of the same trial vectors to $\Xi$ on
  the required two-dimensional compact set;
  \item completion of the infinite-measure Stieltjes compactness and
  one-point criterion in Lean;
  \item independent mathematical review of the analytic bridges and all
  normalization conventions.
\end{enumerate}

Numerical agreement, a finite simple spectrum, or a green Lean build cannot
replace these obligations.  The present result should therefore be read as a
verified finite lemma and a sharper research interface.

\section{Conclusion}

We proved and formalized the finite estimate
\[
 \norm{\norm{k}g_N-k}
 \leq \norm{k}\sqrt{2r_N/\Delta_N}+2\,\mathrm{tail}_N.
\]
The formal statement prevents a common degeneracy error: its positive gap
forces the ground eigenspace to have dimension one.  Because the finite
eigenvalue list is antitone, the full gap premise reduces to one certified
inequality between the last two eigenvalues.  This is a genuine simplification
of the finite verification task.

It is not a proof of the Riemann hypothesis.  The asymptotic same-approximant
convergence and cancellation problem remains the central open step.

\section*{Software, data and reproducibility}

The Lean sources and programme documentation are maintained at
\url{https://github.com/lluiseriksson/riemann-prime-resolvent}.  The campaign
design and witness scripts are maintained in the associated
\path{riemann-resolvent-programme} repository.  Exact source revisions and
registered verdict logs should accompany any archival release of this paper.

\section*{AI-assistance disclosure}

The author directed an AI-assisted drafting and formalization workflow.
Mathematical claims included as Lean theorems were checked by the Lean kernel;
the author remains responsible for the manuscript, its submission, and all
unformalized analytic and numerical statements.

\section*{License}

This manuscript is released under the Creative Commons Attribution 4.0
International license (CC BY 4.0).

\clearpage
\begin{thebibliography}{99}

\bibitem{riemann1859}
B. Riemann,
\emph{Ueber die Anzahl der Primzahlen unter einer gegebenen Gr\"osse},
Monatsberichte der Berliner Akademie, 1859.

\bibitem{titchmarsh1986}
E. C. Titchmarsh and D. R. Heath-Brown,
\emph{The Theory of the Riemann Zeta-Function}, 2nd ed.,
Oxford University Press, 1986.

\bibitem{hausdorff1921}
F. Hausdorff,
Summationsmethoden und Momentfolgen. I,
\emph{Mathematische Zeitschrift} \textbf{9} (1921), 74--109.

\bibitem{shohat1943}
J. A. Shohat and J. D. Tamarkin,
\emph{The Problem of Moments},
American Mathematical Society, 1943.

\bibitem{akhiezer1965}
N. I. Akhiezer,
\emph{The Classical Moment Problem},
Oliver and Boyd, 1965.

\bibitem{kato1995}
T. Kato,
\emph{Perturbation Theory for Linear Operators},
Springer, 1995 reprint of the 1980 edition.

\bibitem{demoura2015}
L. de Moura, S. Kong, J. Avigad, F. van Doorn and J. von Raumer,
The Lean theorem prover (system description),
in \emph{CADE-25}, Lecture Notes in Computer Science 9195,
Springer, 2015, 378--388.

\bibitem{mathlib2020}
The Mathlib Community,
The Lean mathematical library,
in \emph{Proceedings of CPP 2020}, ACM, 2020, 367--381.

\bibitem{connes2025}
A. Connes, C. Consani and H. Moscovici,
Zeta spectral triples,
arXiv:2511.22755, 2025.

\end{thebibliography}

\end{document}
